\documentclass[11pt,a4paper]{article}
\usepackage[margin=28mm]{geometry}
\usepackage[T1]{fontenc}
\usepackage{lmodern}
\pdfextension mapfile{+lm.map}
\pdfextension mapfile{+cm.map}
\pdfextension mapfile{+symbols.map}
\pdfextension mapfile{+rsfs.map}
\pdfextension mapfile{+euler.map}
\usepackage{amsmath,amssymb,amsthm,mathtools,mathrsfs}
\usepackage{booktabs,longtable,array,microtype,enumitem}
\usepackage[hidelinks]{hyperref}
\usepackage{bookmark}
\setlength{\emergencystretch}{3em}
\newenvironment{statement}[3]{\par\medskip\noindent\textbf{Statement #1 (#2). #3.}\ \itshape}{\par\medskip}
\newcommand{\Z}{\mathbb Z}
\newcommand{\Q}{\mathbb Q}
\newcommand{\F}{\mathbb F}
\newcommand{\G}{\operatorname{G}}
\newcommand{\Sec}{\operatorname{Sec}}
\newcommand{\im}{\operatorname{im}}
\newcommand{\coker}{\operatorname{coker}}
\newcommand{\cts}{\mathrm{cts}}
\newcommand{\Frob}{\operatorname{Frob}}
\newcommand{\Rlim}{\mathrm{R}\!\varprojlim}
\newcommand{\limone}{\varprojlim{}^{1}}
\newcommand{\Kcal}{\mathcal K}
\newcommand{\Rcal}{\mathcal R}
\title{LM7 PUSH: arithmetic sections and profinite images\\\large Research manuscript with an explicit statement register}
\author{AI-produced by GPT-6 Astra (OpenAI, ChatGPT Pro mode)}
\date{13 November 2026}
\begin{document}
\maketitle
\begin{abstract}
We compute the intrinsic class-two section obstruction for the Fermat cubic with its nine coordinate cusps removed, including its integral residue complex, path origins, restricted convergence, and lifting fibers. We construct an actual obstruction invisible locally and at every finite coefficient level; refute the two proposed class-two inertia restrictions; and prove the pro-three cyclotomic centralizer criterion. Finite arithmetic field and cocycle constructions give the exact full, metabelian and marked-depth images and all their fibers, with compatibility for arbitrary ramification support. An explicit section at seventeen has the literal local cusp restriction and a global class-three lift, but no class-four lift. We identify the inherited quotient maps and every marked coefficient layer integrally. Over every number field we give the complete arithmetic obstruction test, restriction and descent maps, and the class-two germ classification, including finite algebraic formulas for every point's central coordinates. The three requested full local profinite fixed groups are computed through an explicit valuation and stabilizer comparison. Parts A, B, C and F are complete. The remaining full profinite image questions in Parts D and E are under investigation. We compute the elementary abelian peripheral lifting obstruction, its exact covering restriction, and coupled finite central counterexamples. A transitive $A_5$ wreath family has generating lifts with prescribed peripheral conjugators. The joint binary and ternary degree-six braid refinement retains all twenty-four CM candidates and $864$ finite points. All $432$ multiplier-one points lift through the entire Fermat boundary homology tower, retaining their finite pentagons and generating maps. This working manuscript is \textbf{INCOMPLETE}; the task register identifies every unresolved item. We claim no priority.
\end{abstract}
\tableofcontents
\section*{Conventions and scope}
All cohomology of profinite groups with finite discrete coefficients is continuous. For a free finite-rank $\Z_\ell$-module with continuous action, $H^i_{\cts}$ is the cohomology of continuous inhomogeneous cochains. Restriction, inflation, quotient, and comparison maps are identified in each statement. Section sets are taken modulo conjugation by the specified geometric kernel. Arithmetic Frobenius is used in residue complexes.
\medskip
We work with
\[
 E:\ x^3+y^3+z^3=0,\qquad U=E\setminus D,
\]
where $D$ is the set of nine coordinate cusps. Put
$K_0=\Q(\zeta)$, $\zeta^2+\zeta+1=0$, and $\lambda=1-\zeta$.
Our cusp labels are
\[
\begin{array}{c|ccc}
 &\text{rational cusp}&\text{first conjugate}&\text{second conjugate}\\\hline
 x=0&P=[0:-1:1]&A=[0:-\zeta:1]&\bar A=[0:-\zeta^2:1]\\
 y=0&Q_c=[-1:0:1]&B=[-\zeta:0:1]&\bar B=[-\zeta^2:0:1]\\
 z=0&O=[1:-1:0]&C=[1:-\zeta:0]&\bar C=[1:-\zeta^2:0].
\end{array}
\]
The geometric group $N$ is the free pro-three group of rank ten
of $U$, based at the unit tangent $s=z/x$ at $O$.
Its abelianization and the proper elliptic quotient are
$V=T_3J_{E,D}$ and $T=T_3E$. The cusp submodule is
$W=\Z_3(1)[D]/\Z_3(1)\delta$, with
$\delta=\sum_{d\in D}d$; its inertia vectors are $e_d$.
Let $I=\ker(N\to T)$ and
\[
 N^{\rm cc}=N/[N,I],\qquad
 \mathcal K=\ker\bigl(\bigwedge^2V\to\bigwedge^2T\bigr).
\]
Thus $\mathcal K$ has rank $44$ and is the central kernel
of $N/\gamma_3N\to N^{\rm cc}$. The marked subgroup
$I_{\rm mark}$ is the closed normal subgroup generated by
inertia at $Q_c$. A section class over $L$ is the continuous
nonabelian $H^1$ of $G_L$ with the displayed geometric quotient.
The full cyclotomic character, with target $\widehat\Z^\times$,
is used in the centralizer statements; its three-adic component
acts on the coefficient lattices.

The inherited group $H$ has underlying set $\Z_3^3$ and law
\[
 (u,v,w)(u',v',w')=
 (u+u',v+v',w+w'+\tfrac92(uv'-vu')),
\]
with arithmetic weights $(\chi_3,\chi_3,\chi_3^2)$.
Its base-origin comparison is made explicitly in B.9.
The use of the cuspidally central lifting theorem over $\Q$
is Borne--Emsalem--Stix, Theorem 5.10
\cite{BES2015}, applied to the nine torsion cusps and followed
by the displayed pro-three quotient. Its proof uses vanishing
of relative Brauer groups in all section neighbourhoods.
The point $O$ gives a retraction on the Brauer group of $E$
itself; it does not supply those retractions for every such
neighbourhood. Part C consequently retains the additional
arithmetic terms over a general number field.
For the profinite problems, $\widehat F_2$ is the full free
profinite group on two peripheral generators; the special
outer group and its $\mathfrak S_3$ centralizer
$\widehat{\rm GT}_0$ use all finite quotients, with the latter
defined by the two hexagons. The subgroup $\widehat{\rm GT}$
also satisfies the pentagon. Parts D and E specify every
comparison to a pro-$\ell$, graded, or finite quotient.

The background labels ``LM7'', ``CK'', and ``F17'' refer to the
attached previous returns, and are checked only at the steps
actually used. The available manuscript source is the supplied
transcription of \emph{La Longue Marche}, cote 140-4.
The authoritative scans were not attached. Our manuscript
citations identify their section, folios, scans, and formulas
without asserting a fresh comparison to the scans.
\section{A. The intrinsic class-two problem}
\InputIfFileExists{sections/A_agent.tex}{}{}
\InputIfFileExists{sections/A_root.tex}{}{}
\section{B. Geometric classes and centralizers}
\InputIfFileExists{sections/B_agent.tex}{}{}
\section{C. Restriction and germs}
\InputIfFileExists{sections/C_agent.tex}{}{}
\InputIfFileExists{sections/C_root.tex}{}{}
\InputIfFileExists{sections/C_arithmetic_root.tex}{}{}
\section{D. Profinite leading terms}
\InputIfFileExists{sections/D_agent.tex}{}{}
\InputIfFileExists{sections/D_root.tex}{}{}
\InputIfFileExists{sections/D_peripheral_agent.tex}{}{}
\InputIfFileExists{sections/D_cycle_agent.tex}{}{}
\InputIfFileExists{sections/D_covers_root.tex}{}{}
\InputIfFileExists{sections/D_coupled_agent.tex}{}{}
\InputIfFileExists{sections/D_wreath_agent.tex}{}{}
\section{E. The complex-multiplication coset}
\InputIfFileExists{sections/E_agent.tex}{}{}
\InputIfFileExists{sections/E_root.tex}{}{}
\InputIfFileExists{sections/E_ferrari_degree4_agent.tex}{}{}
\InputIfFileExists{sections/E_ferrari_degree5_agent.tex}{}{}
\InputIfFileExists{sections/E_ferrari_degree6_agent.tex}{}{}
\InputIfFileExists{sections/E_ferrari_two_agent.tex}{}{}
\InputIfFileExists{sections/E_geometry_agent.tex}{}{}
\InputIfFileExists{sections/E_ferrari_s3_agent.tex}{}{}
\InputIfFileExists{sections/E_boundary_agent.tex}{}{}
\InputIfFileExists{sections/E_ferrari_a5_agent.tex}{}{}
\InputIfFileExists{sections/E_torsion_root.tex}{}{}
\InputIfFileExists{sections/E_odd_boundary_agent.tex}{}{}
\InputIfFileExists{sections/E_ferrari_torsion_agent.tex}{}{}
\section{F. Derived limits and applications}
\InputIfFileExists{sections/F_root.tex}{}{}
\InputIfFileExists{sections/F3_action.tex}{}{}
\InputIfFileExists{sections/F3_root.tex}{}{}
\InputIfFileExists{sections/F4_agent.tex}{}{}
\appendix
\section{Coverage and remaining questions}
\InputIfFileExists{sections/coverage.tex}{}{}
\bibliographystyle{alpha}
\bibliography{LM7_PUSH_references}
\end{document}
